\documentclass[12pt]{article}
\usepackage[utf8]{inputenc}
\usepackage[spanish]{babel}
\usepackage{graphicx}
\usepackage{float}
\usepackage{caption}
\usepackage[backend=biber,style=apa]{biblatex}
\addbibresource{../../referencias/referencias.bib}
\usepackage{geometry}
\usepackage{url}
\geometry{margin=2.5cm}



\begin{document}
\begin{titlepage}
\centering
\vspace*{\fill}

{\Huge\bfseries Bitácora 2\par}
\vspace{1.5cm}

{\Large\textbf{Integrantes:}\par}
Debbie Con Ortega (C32250)\\
Ashly Garro Villanueva (C33198)\\
Emily Sanchéz Mancía (C27260)\\
Alessandro Umaña Vega (C37963)

\vspace{1cm}

{\Large\textbf{Profesor:} Maikol Solís\par}
\vspace{1.5cm}

{\Large I Ciclo 2026\par}

\vspace*{\fill}
\end{titlepage}
\newpage
\textbf{Audiencia del proyecto:} \\
\textbf{Respuesta a retroalimentación:} 
\section{Parte de planificación}
\subsection{Fichas de literatura}
\textbf{Ficha 1:}

\textbf{Nivel descriptivo}
\begin{itemize}
    \item \textbf{Título:} Traffic Accidents Analysis on Dry and Wet Road Bends Surfaces in Greater Manchester-UK. 
    
    \item \textbf{Autores:} Ismael Karzan S, Razzaq Noorance A. 
    
    \item \textbf{Año:} 2017 
    
    \item \textbf{Nombre del tema:} Accidentes de tránsito en curvas de Greater Manchester en función de condiciones 
    climáticas y factores de entorno como iluminación, velocidad y comportamiento humano. 
    
    \item \textbf{Clasificación:} Temático. El artículo hace uso de herramientas estadísticas (como la prueba Chi-Cuadrado) y 
    análisis descriptivo para estudiar la relación entre distintas variables, lo cual es directamente aplicable al enfoque 
    de nuestro proyecto. A pesar de esto, se consideran más importantes sus conclusiones para identificar los factores que 
    afectan en mayor magnitud la frecuencia de accidentes de tránsito, por lo que se clasifica como temático.
    
    \item \textbf{Resumen en una oración:} Analiza cómo condiciones de carretera afectan accidentes bajo distintas condiciones 
    climáticas y de entorno. 
    
    \item \textbf{Argumento central:} Las condiciones de la superficie vial (seca o húmeda), junto a otros factores como la 
    iluminación, velocidad y las decisiones del conductor, influyen significativamente en la ocurrencia de accidentes, y estas 
    relaciones pueden evaluarse mediante análisis descriptivo y métodos estadísticos. 
    
    \item \textbf{Problemas con el argumento o el tema:} Se limita a una región específica (Greater Manchester), lo que reduce 
    la generalización de las conclusiones. Por otra parte, considera condiciones climáticas limitadas, dejando por fuera 
    situaciones de alto riesgo como la nevisca. Además, algunas relaciones, como la superficie vial y la severidad de los 
    accidentes, no resultan estadísticamente significativas lo que limita conclusiones fuertes.  \\
\end{itemize}

\textbf{Nivel analítico}

\begin{itemize}
    \item \textbf{Conexión con mi proyecto:} El artículo es sumamente relevante ya que analiza, mediante pruebas estadísticas 
    aplicables para nuestro proyecto, variables similares a las consideradas para nuestra investigación, como condiciones 
    climáticas e iluminación, usando datos reales del Reino Unido (STATS19). 
    
	\item \textbf{Lo que NO dice el autor:} A diferencia del artículo, nuestro proyecto busca explicar cómo varía la cantidad
    total de accidentes en el Reino Unido (no solo en curvas de Greater Manchester), utilizando la base de datos UK Road 
    Safety: Traffic Accidents and Vehicles. Además, no solo describiremos factores asociados (como clima o condición de la 
    carretera), sino que buscaremos identificar situaciones de riesgo recurrentes para proponer medidas de prevención. 
    
	\item \textbf{Contraste con otra fuente:} El estudio “Accident risk of road and weather conditions on different road 
    types” por Malin et al. analiza el riesgo relativo de accidentes considerando múltiples condiciones climáticas (nieve, 
    hielo, visibilidad, temperatura) y distintos tipos de carretera, incorporando además la exposición real al riesgo mediante 
    probabilidad de Palm. En contraste, el estudio de Greater Manchester se enfoca en identificar asociaciones entre variables
    categóricas (superficie seca y húmeda, iluminación, velocidad) mediante pruebas Chi-Cuadrado en un contexto local. 
    Sin embargo, ambos coinciden en que las condiciones de la vía y el entorno influyen significativamente en la ocurrencia 
    de accidentes, y en que el análisis estadístico de estas variables permite identificar factores de riesgo relevantes, 
    aunque con distintos niveles de profundidad.  
    
	\item \textbf{Evaluación de aplicabilidad (1-5):} 4/5. La metodología del artículo es compatible con nuestro proyecto 
    debido a que utiliza pruebas de independencia sobre variables categóricas, que coinciden con la estructura de la base 
    de datos utilizada para nuestro proyecto. Sin embargo, no cumple todos los puntos de aplicabilidad debido a que el 
    enfoque del autor es local y específico a curvas en Greater Manchester mientras que nuestro proyecto abarca todo el 
    Reino Unido y un conjunto de factores más amplio, lo que exige adaptar y ampliar el diseño analítico del artículo.
    
	\item \textbf{Pregunta que le haría al autor:} Dado que su análisis se centra en la relación de la superficie vial en 
    curvas de Greater Manchester tomando en cuenta el factor humano para estudiar las causas de estos, ¿cómo ajustaría su 
    metodología si se quisiera estudiar solamente factores ambientales e infraestructurales, excluyendo el comportamiento 
    del conductor, y qué variables consideraría más relevantes para poder aplicar su enfoque a una base más amplia como la 
    utilizada en nuestro proyecto?
    
	\item \textbf{Resumen argumentativo:} Los autores buscan los factores que influyen en la ocurrencia de accidentes de 
    tránsito en Greater Manchester, particularmente en curvas, analizando asociaciones entre variables como la superficie 
    de la carretera, iluminación, velocidad y severidad del accidente, utilizando pruebas de independencia como la Chi-Cuadrado.
    La solución propuesta resulta parcialmente aplicable a nuestro contexto, ya que el enfoque metodológico basado en el
    análisis de variables categóricas y la identificación de asociaciones entre factores ambientales es coherente con los 
    objetivos de nuestro proyecto. Sin embargo, presenta limitaciones importantes, dado que el estudio se restringe a curvas 
    de Greater Manchester, y considera el comportamiento humano como un factor central, mientras que en nuestra investigación
    se centra en variables climáticas y del entorno, analizando accidentes en todo el Reino Unido. En nuestro trabajo 
    retomaremos el enfoque de análisis asociativo entre variables, por ejemplo, entre condiciones climáticas y cantidad 
    de accidentes, así como la identificación de factores de riesgo, sin embargo, lo aplicaremos a un conjunto de datos 
    más amplio y sin atribuir la causalidad principal al error del conductor.\\
\end{itemize} 

\textbf{Ficha 5:}

\textbf{Nivel descriptivo}
\begin{itemize}
    \item \textbf{Título:} Análisis estadístico de tablas de contingencia y chi-cuadrado para medir el flujo migratorio de
    origen y destino en el Ecuador año 2018.
    
    \item \textbf{Autores:}  David Lastre Baquerizo, Melanie Páez Santana, y Olga López Tumbaco.
    
    \item \textbf{Año:} 2019
    
    \item \textbf{Nombre del tema:} Aplicación de tablas de contingencia y prueba chi-cuadrado para analizar el flujo 
    migratorio en Ecuador.
    
    \item \textbf{Clasificación:} Metodológico. El artículo explica detalladamente el funcionamiento de las tablas de contingencia y de 
    la prueba de independencia Chi-Cuadrado.
    
    \item \textbf{Resumen en una oración:} Uso de pruebas de independencia Chi-cuadrado para analizar patrones de flujo migratorio en Ecuador.
    
    \item \textbf{Argumento central:} El artículo sostiene que el uso de herramientas estadísticas como las tablas de contingencia y la prueba 
    chi-cuadrado permite analizar de manera objetiva y eficiente los patrones de migración entre origen y destino, facilitando la interpretación 
    de relaciones entre variables y apoyando la toma de decisiones basada en datos.
    
    \item \textbf{Problemas con el argumento o el tema:} El estudio muestra limitaciones relacionadas con la disponibilidad y calidad de los datos 
    históricos, además de la dependencia de supuestos estadísticos que pueden afectar la precisión del análisis. También puede existir dificultad para 
    explicar la naturaleza social del fenómeno migratorio únicamente mediante métodos cuantitativos.\\
\end{itemize}

\textbf{Nivel analítico}

\begin{itemize}
    \item \textbf{Conexión con mi proyecto:} El artículo se conecta con nuestro proyecto en la metodología que utiliza, estas son tablas de contingencia 
    y pruebas de Chi-Cuadrado para encontrar patrones de migración (de salida o entrada) por sexo, nacionalidad, motivo de viaje, entre otras. 
    
	\item \textbf{Lo que NO dice el autor:} El artículo no responde a ningún aspecto de la pregunta de investigación, ni permite identificar relaciones 
    específicas entre variables propias de este fenómeno. Su principal limitación para este proyecto es que, aunque desarrolla herramientas metodológicas como 
    las tablas de contingencia y la prueba chi-cuadrado, no las aplica en el contexto de accidentes de tránsito. 
    
	\item \textbf{Contraste con otra fuente:} El artículo de Lastre et al. propone el uso de pruebas de independencia Chi-Cuadrado para analizar la relación 
    entre variables categóricas en un contexto de migración. Por otro lado, el artículo de Montaña et al. (2025) se enfoca en la prueba de homogeneidad 
    Chi-Cuadrado, en el ámbito de estudios clínicos. Ambos documentos se complementan ya que abordan distintas aplicaciones de la misma familia de métodos 
    estadísticos. El primero se centra en encontrar una relación significativa entre variables, es decir, cuantificar la dependencia o independencia de las 
    variables, mientras que el segundo se enfoca en determinar si existen diferencias o similitudes en la distribución de distintas categorías de las variables. 
    Esto es relevante para nuestro proyecto pues no solo nos interesa encontrar relación entre distintos factores y la frecuencia de accidentes de tránsito, sino 
    también comparar cómo varían y se distribuyen los accidentes entre distintos grupos y condiciones. Además, el segundo artículo ayuda a resolver una limitación 
    del primero, ya que amplía el enfoque metodológico al introducir la prueba de homogeneidad como herramienta adicional para el análisis de variables categóricas. 
    Sin embargo, ambos comparten el problema de no abordar directamente el tema de accidentes de tránsito, por lo que su aporte se mantiene en el plano metodológico.
    
	\item \textbf{Evaluación de aplicabilidad (1-5):} La metodología es aplicable porque tanto el artículo como el proyecto trabajan con variables categóricas, lo que 
    hace pertinente el uso de tablas de contingencia y la prueba chi-cuadrado.  Sin embargo, no se asigna un 5 porque el artículo no aborda las particularidades de los 
    datos de accidentes de tránsito. 
    
	\item \textbf{Pregunta que le haría al autor:} ¿Cómo recomienda adaptar la prueba chi-cuadrado cuando se trabaja con datos de accidentes de tránsito que pueden presentar 
    baja frecuencia en ciertas categorías o posible dependencia entre observaciones?
    
	\item \textbf{Resumen argumentativo:} El autor busca resolver el problema de cómo analizar la relación entre variables categóricas mediante herramientas estadísticas 
    como las tablas de contingencia y la prueba chi-cuadrado para explicar el flujo migratorio en Ecuador. Su propuesta funciona porque estas técnicas permiten cuantificar 
    la asociación entre variables, permitiendo interpretar patrones dentro de los datos. En nuestro proyecto, la solución es adecuada en cierta medida, ya que las variables 
    con las que se trabajarán también son categóricas y el objetivo es identificar relaciones entre factores ambientales y la frecuencia de accidentes. Sin embargo, el enfoque 
    del artículo no considera la complejidad propia de los accidentes de tránsito, como su naturaleza multicausal o la presencia de dependencias entre observaciones. A pesar de 
    esto, la metodología aplicada en el artículo sigue siendo aceptable como herramienta de análisis. Específicamente, el uso de tablas de contingencia permitirá organizar y 
    visualizar las frecuencias cruzadas entre variables importantes. Asimismo, la prueba chi-cuadrado será utilizada para determinar si existe evidencia estadística de asociación 
    entre estas variables. Por lo tanto, el principal aporte del artículo al proyecto es proporcionar una base metodológica completa para el análisis de nuestras variables.\\
\end{itemize} 

\subsection{Ordenamiento de la literatura}
\begin{table}[H]
\centering
\renewcommand{\arraystretch}{1.5}

\begin{tabular}{|p{2.5cm}|p{5cm}|p{4.5cm}|p{1cm}|p{3cm}|}
\hline
\textbf{Clasificación} & \textbf{Justificación breve} & \textbf{Título} & \textbf{Año} & \textbf{Autor(es)} \\
\hline
Temático & Análisis estadístico (Chi-Cuadrado) de variables categóricas como clima, iluminación y 
superficie vial sobre accidentes en curvas & Traffic Accidents Analysis on Dry and Wet Road Bends Surfaces in Greater 
Manchester-UK & 2017 & Ismael Karzan S., Razzaq Noorance A. \\
\hline
Temático & Extensión del enfoque anterior: cuantifica el riesgo relativo según condiciones climáticas y tipo de 
carretera usando distribución de Palm & Accident risk of road and weather conditions on different road types & 2019 & 
Fanny Malin, Ilkka Norros, Saty Innamaa \\
\hline
Teórica & Marco conceptual para justificar la velocidad como factor en la frecuencia de accidentes; sustenta hipótesis 
sobre tipos de carretera y zonas urbanas/rurales & The effects of drivers' speed on the frequency of road accidents & 2000 
& M. C. Taylor, D. A. Lynam, A. Baruya \\
\hline
Teórica & Marco conceptual para justificar el análisis multifactorial de accidentes mediante la Teoría del Dominó y la Homeostasis
 del Riesgo & Analysis of Road Traffic Accident using Causation Theory with Traffic Safety Model and Measures & 2017 & Dr. Neeta Saxena \\
\hline
\end{tabular}
\end{table}

\section{Enlaces de la literatura}


\section{Análisis estadísticos}
\subsection{Indentidad visual del proyecto}
En este proyecto se definió un tema para la visualización de datos denominado estilo_bayesianos, con el objetivo de asegurar coherencia y claridad en todos los gráficos. Para ello, se adoptó un diseño minimalista que elimina elementos innecesarios, como las grillas de fondo, y prioriza la legibilidad de la información mediante ejes bien definidos y tamaños de fuente adecuados para su correcta visualización en el documento final. Asimismo, la leyenda se ubica en la parte superior para facilitar la interpretación de los gráficos sin interferir con el área de datos. En cuanto a la selección de colores, se emplea la paleta BMJ del paquete ggsci, la cual se caracteriza por ofrecer colores contrastantes, armónicos y fácilmente distinguibles. Esta elección resulta especialmente adecuada para representar variables categóricas, ya que mejora la diferenciación visual entre grupos y contribuye a una presentación profesional y consistente de los resultados. \\

\subsection{Análisis descriptivo}
\subsection{Construcción de fichas de resultados}
\subsection{La imagen que cuenta la historia}


\section{Arco narrativo del proyecto}

\section{Diario de aprendizaje}

\newpage
\printbibliography


\newpage
\subsection*{Log de contribución}

\begin{table}[h!]
\centering
\small
\begin{tabular}{|p{3cm}|p{8cm}|}
\hline
\textbf{Integrante} & \textbf{Tareas realizadas} \\ \hline

Debbie Con & 
Definición de la idea \par
Reformulación de la idea en modo de pregunta \par
Argumentación de pregunta 4 \par
Ficha literaria 3 \par
Construcción de la UVE de Gowin \par
Redacción del resumen \par
Arco narrativo del proyecto \par
Registro del uso de IA \par
Identificación de tensiones \\ \hline

Ashly Garro & 
Definición de la idea \par
Reformulación de la idea en modo de pregunta \par
Argumentación de pregunta 2 \par
Ficha literaria 1 \par
Construcción de la UVE de Gowin \par
Arco narrativo del proyecto \par
Rastro de decisiones \par
Diario de aprendizaje \\ \hline

Emily Sánchez &
Definición de la idea \par
Reformulación de la idea en modo de pregunta \par
Conceptualización de la idea \par
Argumentación de la pregunta 3 \par
Construcción de la UVE de Gowin \par
Arco narrativo del proyecto \par
Ficha literaria 2 \par
Diario de aprendizaje \\ \hline

Alessandro Umaña & 
Argumentación de la pregunta 1 \par
Argumentación a través de los datos \par
Construcción de la UVE de Gowin \par
Ficha literaria 4 \par
Arco narrativo del proyecto \par
Diario de aprendizaje \\ \hline

\end{tabular}
\end{table}

\subsection*{Rastro de decisiones}

\subsection*{Registro de uso de IA}

\textbf{Primera pregunta}
\begin{itemize}
    \item Herramienta:  
    \item Prompt: 
    \item Respuesta de la IA: 
    \item Evaluación: 
    \item Modificaciones: 
    \item Aprendizaje:
\end{itemize}

\textbf{Segunda pregunta:}
\begin{itemize}
    \item Herramienta: 
    \item Prompt: 
    \item Respuesta de la IA: 
 
    \item Evaluación: 
    \item Modificaciones: 
    \item Aprendizaje:
\end{itemize}
\end{document}